\pagestyle{empty}
\tight
\begin{tblr}{width=\textwidth,colspec={|Q[c,m,1.9cm]|X[l,m]|},hlines,vlines,hspan=minimal,row{1}={bg=headblue},rowsep=1pt,colsep=3pt}
\SetCell{c,m}\hfil\logo\hfil & \textbf{POLITEKNIK NEGERI MALANG}\par\textbf{JURUSAN TEKNOLOGI INFORMASI}\par\textbf{PROGRAM STUDI: D4 TEKNIK INFORMATIKA} \\
\SetCell[c=2]{c} \textbf{RENCANA PEMBELAJARAN SEMESTER (RPS)} & \\
\end{tblr}
\begin{longtblr}{width=\textwidth,colspec={|Q[l,m,1.9cm]|X[0.85,l,m]|X[1.45,l,m]|X[0.75,l,m]|X[1.15,l,m]|},hlines,vlines,hspan=minimal,rowsep=1pt,colsep=2pt,row{1,3}={bg=headgray,font=\bfseries}}
MATA KULIAH & KODE & BOBOT (SKS)/jam & SEMESTER & TGL. PENYUSUNAN \\
Praktikum Jaringan Komputer & RTI254006 & 2 SKS/4 jam & 4 & 1 Juli 2025 \\
OTORISASI & Dosen Pengembang RPS & Koordinator RMK & \SetCell[c=2]{l} Ka PRODI &  \\
 & Vipkas Al Hadid Firdaus, S.T., M.T. & Vipkas Al Hadid Firdaus, S.T., M.T. & \SetCell[c=2]{l}{\vspace{8mm}\par Dr. Ely Setyo Astuti, S.T., M.T.} &  \\
\SetCell[r=9]{l} Capaian Pembelajaran (CP) & \SetCell[c=4]{l,bg=headgray,font=\bfseries} Capaian Pembelajaran Lulusan Program Studi (CPL-Prodi) &  &  &  \\
 & CPL05 & \SetCell[c=3]{l} Mampu menyediakan solusi teknologi komputasi dan infrastruktur digital yang memenuhi kebutuhan pengguna dan organisasi. &  &  \\
 & CPL09 & \SetCell[c=3]{l} Mampu memahami cara kerja sistem komputer dan menerapkan pendekatan komputasi untuk problem solving. &  &  \\
 & \SetCell[c=4]{l,bg=headgray,font=\bfseries} Capaian Pembelajaran mata kuliah (CPL-MK) &  &  &  \\
 & CPMK0504 & \SetCell[c=3]{l} Mahasiswa mampu memilih arsitektur, protokol, dan mengkonfigurasi infrastruktur jaringan komputer secara praktis. &  &  \\
 & CPMK0901 & \SetCell[c=3]{l} Mahasiswa mampu menerapkan cara kerja sistem jaringan komputer melalui praktikum konfigurasi dan troubleshooting. &  &  \\
 & \SetCell[c=4]{l,bg=headgray,font=\bfseries} Kemampuan akhir tiap tahapan belajar (Sub-CPMK) &  &  &  \\
 & SCPMK0504-03301 & \SetCell[c=3]{l} Mahasiswa mampu mengkonfigurasi perangkat jaringan (NIC, router, switch) dan melakukan pengujian konektivitas. &  &  \\
 & SCPMK0504-03302 & \SetCell[c=3]{l} Mahasiswa mampu melakukan subnetting IPv4 dan konfigurasi routing dasar menggunakan Packet Tracer. &  &  \\
 & SCPMK0901-03303 & \SetCell[c=3]{l} Mahasiswa mampu menganalisis traffic jaringan menggunakan tools (ping, traceroute, ARP) dan mendokumentasikan hasilnya. &  &  \\
\SetCell[r=2]{l} Penilaian & \SetCell[c=4]{l} *Nilai CPMK didapat dari agregasi nilai SUB-CPMK yang dibebankan &  &  &  \\
 & \SetCell[c=4]{l}{\tiny\begin{tblr}{width=\linewidth,colspec={|Q[c,m,0.09\linewidth]|Q[c,m,0.15\linewidth]|X[l,m]|X[l,m]|X[l,m]|X[l,m]|X[l,m]|X[l,m]|Q[c,m,0.08\linewidth]|},hlines,vlines,hspan=minimal,rowsep=1pt,colsep=0.7pt,row{1,2}={bg=headgray,font=\bfseries}}
Kode CPMK & Kode SCPMK & \SetCell[c=6]{c} Bobot per Bentuk Penilaian & & & & & & Total Bobot \\
 & & Tugas & Kuis 1 & UTS & Kuis 2 & PBL & UAS & \\
\SetCell[r=2]{c} CPMK0504 & SCPMK0504-03301 & 7\%: laporan konfigurasi media jaringan & & & & 18\%: jobsheet konfigurasi NIC, switch, dan konektivitas & & 25\% \\
 & SCPMK0504-03302 & 7\%: laporan subnetting dan routing & 5\%: kuis konsep jaringan & & & 19\%: jobsheet subnetting IPv4 dan Packet Tracer & & 31\% \\
\SetCell[r=1]{c} CPMK0901 & SCPMK0901-03303 & 6\%: laporan analisis traffic jaringan & & & 5\%: kuis troubleshooting jaringan & 18\%: jobsheet analisis ping, traceroute, ARP & 15\%: ujian akhir praktikum & 44\% \\
\SetCell[c=8]{r}\textbf{Total Per Penilaian} & & & & & & & & \textbf{100\%} \\
\end{tblr}} &  &  &  \\
Diskripsi Singkat Mata Kuliah & \SetCell[c=4]{l} Praktikum Jaringan Komputer memberikan pengalaman langsung dalam mengkonfigurasi media jaringan, protokol TCP/IP, subnetting IPv4, routing, dan simulasi jaringan menggunakan Cisco Packet Tracer. &  &  &  \\
Bahan Kajian / Materi Pembelajaran & \SetCell[c=4]{l}\begin{enumerate}\item Media dan konfigurasi perangkat jaringan\item Protokol lapisan aplikasi dan transport\item IPv4, subnetting, dan routing\item Packet Tracer dan analisis traffic\end{enumerate} &  &  &  \\
\SetCell[r=2]{l} Pustaka & \SetCell[c=4]{l} \textbf{Utama:}\begin{enumerate}\item Cisco Networking Academy. 2020. CCNA: Introduction to Networks. Cisco Press.\item Forouzan, B.A. 2012. Data Communications and Networking. McGraw-Hill.\item Modul praktikum jaringan komputer JTI Polinema.\end{enumerate} &  &  &  \\
 & \SetCell[c=4]{l} \textbf{Pendukung:} Materi LMS, dokumentasi resmi perangkat lunak, dan jobsheet praktikum. &  &  &  \\
Media Pembelajaran & \SetCell[c=2]{l} \textbf{Software:} Cisco Packet Tracer, Wireshark. &  & \SetCell[c=2]{l} \textbf{Hardware:} Komputer lab, perangkat jaringan. &  \\
Nama Dosen Pengampu & \SetCell[c=4]{l} Tim Pengajar Program Studi D4 TI &  &  &  \\
Matakuliah Syarat & \SetCell[c=4]{l} RTI254005 Jaringan Komputer (corequisite). &  &  &  \\
\end{longtblr}
\begin{longtblr}{width=\textwidth,colspec={|Q[c,m,0.06\textwidth]|X[1.35,l,m]|X[1.08,l,m]|X[1.16,l,m]|X[0.7,l,m]|X[1.24,l,m]|X[1.06,l,m]|X[0.98,l,m]|Q[c,m,0.05\textwidth]|},hlines,vlines,hspan=minimal,rowhead=2,row{1,2}={bg=headgreen,font=\bfseries},rowsep=1pt,colsep=1.5pt}
Mgg.\par Ke & Kemampuan Akhir Yang Direncanakan (Sub-CP-MK) & Bahan kajian (Materi Pembelajaran) & Bentuk dan Metode Pembelajaran & Estimasi Waktu & Pengalaman Belajar Mahasiswa & Kriteria \& Bentuk Penilaian & Indikator Penilaian & Bobot Penilaian (\%) \\
(1) & (2) & (3) & (4) & (5) & (6) & (7) & (8) & (9) \\
1 & SCPMK0504-03301 & Media Jaringan (Network Media) & Modalitas: Luring. Bentuk: Praktikum. Metode: Demonstrasi, praktikum jobsheet, troubleshooting, dan dokumentasi. & 1 x 4 x 50' praktikum & Mahasiswa mengerjakan jobsheet praktikum media jaringan dan mendokumentasikan hasil konfigurasi/analisis. & Bentuk penilaian: Media Jaringan. Mengacu rubrik RTM. & Ketepatan konfigurasi; kualitas jobsheet; validasi dan dokumentasi. & 4.58 \\
2 & SCPMK0504-03301 & Konfigurasi Network Card (Network Setting and Configuration) & Modalitas: Luring. Bentuk: Praktikum. Metode: Demonstrasi, praktikum jobsheet, troubleshooting, dan dokumentasi. & 1 x 4 x 50' & Mahasiswa mengerjakan jobsheet praktikum konfigurasi NIC dan mendokumentasikan hasil konfigurasi/analisis. & Bentuk penilaian: Konfigurasi Network Card. Mengacu rubrik RTM. & Ketepatan konfigurasi; kualitas jobsheet; validasi dan dokumentasi. & 4.58 \\
3 & SCPMK0504-03301 & Application Layer Protocol & Modalitas: Luring. Bentuk: Praktikum. Metode: Demonstrasi, praktikum jobsheet, troubleshooting, dan dokumentasi. & 1 x 4 x 50' & Mahasiswa mengerjakan jobsheet praktikum protokol aplikasi dan mendokumentasikan hasil konfigurasi/analisis. & Bentuk penilaian: Application Layer Protocol. Mengacu rubrik RTM. & Ketepatan konfigurasi; kualitas jobsheet; validasi dan dokumentasi. & 4.58 \\
4 & SCPMK0504-03301 & Transport Layer Protocol & Modalitas: Luring. Bentuk: Praktikum. Metode: Demonstrasi, praktikum jobsheet, troubleshooting, dan dokumentasi. & 1 x 4 x 50' & Mahasiswa mengerjakan jobsheet praktikum protokol transport dan mendokumentasikan hasil konfigurasi/analisis. & Bentuk penilaian: Transport Layer Protocol. Mengacu rubrik RTM. & Ketepatan konfigurasi; kualitas jobsheet; validasi dan dokumentasi. & 4.58 \\
5 & SCPMK0504-03301 & Ping and Route & Modalitas: Luring. Bentuk: Praktikum. Metode: Demonstrasi, praktikum jobsheet, troubleshooting, dan dokumentasi. & 1 x 4 x 50' & Mahasiswa mengerjakan jobsheet praktikum ping dan routing dan mendokumentasikan hasil konfigurasi/analisis. & Bentuk penilaian: Ping and Route. Mengacu rubrik RTM. & Ketepatan konfigurasi; kualitas jobsheet; validasi dan dokumentasi. & 4.58 \\
6 & SCPMK0504-03302 & Kuis 1 / Quiz 1 & Modalitas: Luring. Bentuk: Kuis. Metode: Ujian tertulis konsep jaringan. & 1 x 4 x 50' & Mahasiswa mengerjakan kuis tertulis konsep jaringan layer 2-4. & Kuis 1 & Ketepatan jawaban konsep jaringan & 5 \\
7 & SCPMK0504-03302 & IPv4 & Modalitas: Luring. Bentuk: Praktikum. Metode: Demonstrasi, praktikum jobsheet, troubleshooting, dan dokumentasi. & 1 x 4 x 50' & Mahasiswa mengerjakan jobsheet praktikum konsep dan konfigurasi IPv4 dan mendokumentasikan hasil konfigurasi/analisis. & Bentuk penilaian: IPv4. Mengacu rubrik RTM. & Ketepatan konfigurasi; kualitas jobsheet; validasi dan dokumentasi. & 4.58 \\
8 & SCPMK0504-03302 & IPv4 Subnetting & Modalitas: Luring. Bentuk: Praktikum. Metode: Demonstrasi, praktikum jobsheet, troubleshooting, dan dokumentasi. & 1 x 4 x 50' & Mahasiswa mengerjakan jobsheet praktikum subnetting IPv4 dan mendokumentasikan hasil konfigurasi/analisis. & Bentuk penilaian: IPv4 Subnetting. Mengacu rubrik RTM. & Ketepatan konfigurasi; kualitas jobsheet; validasi dan dokumentasi. & 4.58 \\
9 & SCPMK0504-03301, SCPMK0504-03302 & UTS / Middle Test & Modalitas: Luring. Bentuk: Ujian Praktikum Tengah Semester. Metode: Ujian praktikum individual. & 1 x 4 x 50' & Mahasiswa mengerjakan ujian praktikum mid-semester konfigurasi jaringan. & UTS Praktikum & Ketepatan konfigurasi; kualitas laporan ujian & 10 \\
10 & SCPMK0901-03303 & Traceroute and Tracert & Modalitas: Luring. Bentuk: Praktikum. Metode: Demonstrasi, praktikum jobsheet, troubleshooting, dan dokumentasi. & 1 x 4 x 50' & Mahasiswa mengerjakan jobsheet praktikum traceroute dan analisis path dan mendokumentasikan hasil konfigurasi/analisis. & Bentuk penilaian: Traceroute and Tracert. Mengacu rubrik RTM. & Ketepatan konfigurasi; kualitas jobsheet; validasi dan dokumentasi. & 4.58 \\
11 & SCPMK0901-03303 & Address Resolution Protocol & Modalitas: Luring. Bentuk: Praktikum. Metode: Demonstrasi, praktikum jobsheet, troubleshooting, dan dokumentasi. & 1 x 4 x 50' & Mahasiswa mengerjakan jobsheet praktikum ARP dan analisis address resolution dan mendokumentasikan hasil konfigurasi/analisis. & Bentuk penilaian: Address Resolution Protocol. Mengacu rubrik RTM. & Ketepatan konfigurasi; kualitas jobsheet; validasi dan dokumentasi. & 4.58 \\
12 & SCPMK0504-03302, SCPMK0901-03303 & Packet Tracer (simulasi jaringan) & Modalitas: Luring. Bentuk: Praktikum. Metode: Demonstrasi, praktikum jobsheet, troubleshooting, dan dokumentasi. & 1 x 4 x 50' & Mahasiswa mengerjakan jobsheet praktikum simulasi jaringan Packet Tracer dan mendokumentasikan hasil konfigurasi/analisis. & Bentuk penilaian: Packet Tracer. Mengacu rubrik RTM. & Ketepatan konfigurasi; kualitas jobsheet; validasi dan dokumentasi. & 4.58 \\
13 & SCPMK0504-03302, SCPMK0901-03303 & Packet Tracer lanjutan & Modalitas: Luring. Bentuk: Praktikum. Metode: Demonstrasi, praktikum jobsheet, troubleshooting, dan dokumentasi. & 1 x 4 x 50' & Mahasiswa mengerjakan jobsheet praktikum Packet Tracer lanjutan dan mendokumentasikan hasil konfigurasi/analisis. & Bentuk penilaian: Packet Tracer Lanjutan. Mengacu rubrik RTM. & Ketepatan konfigurasi; kualitas jobsheet; validasi dan dokumentasi. & 4.58 \\
14 & SCPMK0901-03303 & Kuis 2 / Quiz 2 & Modalitas: Luring. Bentuk: Kuis. Metode: Ujian tertulis troubleshooting jaringan. & 1 x 4 x 50' & Mahasiswa mengerjakan kuis troubleshooting dan analisis traffic. & Kuis 2 & Ketepatan analisis dan troubleshooting jaringan & 5 \\
15 & SCPMK0504-03301, SCPMK0504-03302, SCPMK0901-03303 & Presentasi (Presentation) -- presentasi hasil proyek konfigurasi & Modalitas: Luring. Bentuk: Presentasi. Metode: Presentasi proyek konfigurasi jaringan dan peer review. & 1 x 4 x 50' & Mahasiswa mempresentasikan proyek konfigurasi jaringan kepada dosen dan teman sejawat. & Presentasi Proyek & Kualitas presentasi; kedalaman analisis; dokumentasi & 10 \\
16 & SCPMK0504-03302 & Konfigurasi Dasar Router 1 (Basic Router Configuration 1) & Modalitas: Luring. Bentuk: Praktikum. Metode: Demonstrasi, praktikum jobsheet, troubleshooting, dan dokumentasi. & 1 x 4 x 50' & Mahasiswa mengerjakan jobsheet praktikum konfigurasi dasar router dan mendokumentasikan hasil konfigurasi/analisis. & Bentuk penilaian: Konfigurasi Dasar Router 1. Mengacu rubrik RTM. & Ketepatan konfigurasi; kualitas jobsheet; validasi dan dokumentasi. & 4.58 \\
17 & SCPMK0504-03301, SCPMK0504-03302, SCPMK0901-03303 & Konfigurasi Dasar Router 2 / UAS & Modalitas: Luring. Bentuk: Ujian Akhir Semester. Metode: Ujian praktikum final. & 1 x 4 x 50' & Mahasiswa mengerjakan ujian akhir praktikum konfigurasi router dan troubleshooting. & UAS Praktikum & Ketepatan konfigurasi; kualitas laporan; pemecahan masalah & 15 \\
\end{longtblr}
